\documentclass[11pt,a4paper]{article}
\usepackage[margin=22mm,headheight=15pt]{geometry}
\usepackage{fontspec}
\IfFontExistsTF{Malgun Gothic}{\setmainfont{Malgun Gothic}}{\setmainfont{Noto Sans CJK KR}}
\usepackage{kotex}
\usepackage{amsmath,amssymb,array,tabularx,longtable,booktabs,enumitem,xcolor,listings,fancyhdr,hyperref,xurl}
\definecolor{navy}{HTML}{17365D}
\definecolor{light}{HTML}{F3F5F7}
\hypersetup{colorlinks=true,linkcolor=navy,urlcolor=navy,pdftitle={HAIC 제출 코드 전체 해설}}
\Urlmuskip=0mu plus 1mu\relax
\setlength{\parindent}{0pt}\setlength{\parskip}{5pt}
\renewcommand{\baselinestretch}{1.13}
\setlist[itemize]{nosep,leftmargin=1.5em}
\setlist[enumerate]{itemsep=2pt,leftmargin=1.7em}
\pagestyle{fancy}\fancyhf{}\fancyhead[L]{\small HAIC 제출 코드 전체 해설}\fancyhead[R]{\small 2026-09-29}\fancyfoot[C]{\thepage}
\lstset{language=Python,basicstyle=\ttfamily\fontsize{8}{9}\selectfont,breaklines=true,breakatwhitespace=false,columns=fullflexible,keepspaces=true,showstringspaces=false,numbers=left,numberstyle=\tiny\color{gray},numbersep=6pt,xleftmargin=16pt,frame=single,rulecolor=\color{gray},backgroundcolor=\color{light},tabsize=4,keywordstyle=\color{navy}\bfseries,commentstyle=\color{gray},stringstyle=\color{teal}}
\newcommand{\code}[1]{{\small\nolinkurl{#1}}}
\newcommand{\clip}{\operatorname{clip}}
\begin{document}
\begin{titlepage}
\vspace*{18mm}
{\Huge\bfseries\color{navy} HAIC 제출 코드\\[4mm]전체 해설}\par
\vspace{8mm}
{\Large 실행 전략 · 계산식 · 상수 · 상태 · 전체 소스}
\vspace{14mm}

\textbf{문서 기준}\par
현재 제출 추천 패키지 \code{submission-fast-completion.zip} 안에 들어 있는 7개 파일을 직접 읽어 작성했다. 작업 폴더의 최신 후보 코드와 섞지 않았다. 코드의 동작 설명과 전체 원문을 제공하며, 연구 과정이나 실험 이력은 다루지 않는다.

\textbf{실제로 선택된 실행 조합}\par
\code{Agent} $\rightarrow$ \code{FastCompletionCoordination(preview_row_repair)}\\
$\rightarrow$ \code{FixedHighSpeedRecoveryV2(impact_clear)}\\
$\rightarrow$ \code{FixedHighSpeedAgent(damping)}\\
$\rightarrow$ 내부 객체 \code{VisionCorridorAgent}

\textbf{읽는 순서}\par
앞부분은 실제 실행되는 경로를 설명한다. 다음에는 패키지에 포함되지만 선택되지 않은 모드와 보조 함수를 설명한다. 마지막 부록에는 모든 파일의 전체 소스를 줄 번호와 함께 수록한다. 줄 번호는 각 원본 파일의 줄 번호다. 긴 줄의 지면상 줄바꿈은 원본 줄을 추가한 것이 아니다.

\vfill
2026년 9월 29일\\
속도 수치는 게임 내부 단위다. 페달 명령의 크기와 실제 가속도는 구별한다.
\end{titlepage}
\renewcommand{\contentsname}{목차}\tableofcontents\newpage

\section{패키지 구성과 문서 범위}
이 패키지는 NumPy를 사용하는 규칙 기반 제어기다. 신경망 가중치, 학습 모델 로딩, 맵 좌표 파일, 경로 최적화기, 외부 통신은 포함되어 있지 않다. 부록의 전체 소스는 ZIP의 텍스트를 그대로 삽입했다. 공백·주석·사용하지 않는 함수도 생략하지 않았다.

\par\noindent\begin{tabularx}{\linewidth}{@{}p{74mm}X@{}}\toprule
파일 & 역할 \\\midrule
\code{agent.py} & 대회 호출 인터페이스와 실제 모드 선택 \\
\code{haic_agent/__init__.py} & 패키지 설명 문자열 \\
\code{haic_agent/pixel_features.py} & 입력 검사, 도로·장애물·속도 추출과 보조 특징 함수 \\
\code{haic_agent/corridor_agent.py} & 중심선 추종, 회피 방향 상태, 하위 속도 제어와 진단 \\
\code{haic_agent/fixed_high_speed_runtime.py} & 조향 보정과 고정 60 목표의 페달 제어 \\
\code{haic_agent/fixed_high_speed_recovery_v2.py} & 급감속 감지와 회피 조향 해제 \\
\code{haic_agent/fast_completion_coordination.py} & 제동 오인 방지, 누락 도로 복원, 최종 목표 속도 \\\bottomrule
\end{tabularx}

물리 엔진이나 환경 래퍼는 ZIP에 포함되지 않는다. 따라서 그 내부 코드를 이 패키지의 기능으로 설명하지 않는다. 이 문서의 ``모든 코드''는 이 제출 ZIP의 일곱 파일 전체를 의미한다.

\section{진입점, 입력과 출력}
\subsection{Agent의 네 메서드}
\code{__init__}은 \code{FastCompletionCoordination("preview_row_repair")} 객체를 하나 만든다. \code{reset(observation)}은 내부 객체의 reset으로 전달한다. \code{act(observation)}은 내부 객체의 최종 행동을 그대로 반환한다. \code{last_step_diagnostics()}는 내부 진단 사전의 복사본을 반환한다. Agent 자체의 추가 판단은 없다.

\subsection{관측 형식과 검사}
\code{current_frame}은 입력을 NumPy \code{float32} 배열로 변환한다. 변환 중 TypeError 또는 ValueError는 설명 문자열을 붙인 ValueError로 바꾼다. 허용 모양은 $(4,84,84)$ 또는 $(84,84)$이다. 네 장일 때 마지막 장을 현재 화면으로 선택한다. 현재 화면의 모든 값이 유한하고 $[0,1]$에 있어야 한다. 다른 크기나 범위는 예외다.

여기서 직접 검사하는 값은 \textbf{선택된 현재 화면}이다. 네 장 전체를 검증하는 것은 아니다. 이전 화면은 이후 속도 평균과 조향 보정에 사용되므로, 모든 입력 프레임이 정상이라는 환경 계약이 중요하다. 선택된 최상위 실행 경로에는 이 예외를 잡아 행동을 대신 반환하는 처리가 없다.

출력은 \code{np.asarray([steer, gas, brake], dtype=np.float32)}인 길이 3 배열이다. 최종 조향은 보통 $[-0.7,0.7]$, 가속은 $[0,0.6]$, 제동은 $[0,0.15]$ 또는 $[0,0.28]$ 범위다. 이들은 조작 명령이며 차량 위치나 실제 가속도를 직접 지정하지 않는다.

\section{화면에서 도로 찾기}
\subsection{밝기와 탐색 범위}
도로 픽셀은 $0.24\le I(y,x)\le0.52$로 분류한다. 화면 가로 기준은 $x=42$이다. $y$는 아래로 갈수록 증가한다. 탐색 행은 가까운 곳부터 $54,50,46,42,38,34,30$이다.

처음 탐색 중심을 42로 놓고, 각 행에서 이전에 찾은 중심의 좌우 17픽셀 안에 있는 도로 픽셀만 선택한다. 최소 4개 픽셀이 있으면 그 $x$좌표들의 평균을 그 행의 중심으로 저장하고, 다음 행의 탐색 중심도 이 평균으로 바꾼다. 4개 미만이면 그 행은 사전에 넣지 않고 이전 탐색 중심을 유지한다.

이는 연속된 도로 경계의 중점을 엄밀히 구하는 방식은 아니다. 탐색 범위에 포함된 도로 밝기 픽셀의 평균이므로, 가림·다른 도로 조각·밝기 오분류가 중심에 영향을 줄 수 있다.

\subsection{두 가지 보간 함수}
\code{center_at(row, centers)}는 알려진 행을 오름차순 정렬하고 \code{np.interp}로 중심을 선형 보간한다. 알려진 행이 없으면 42를 반환한다. 범위 밖에서는 가장 끝의 값으로 유지된다. 장애물 높이에서 도로 중심을 구할 때 사용한다.

\code{observed_center(centers,row)}는 요청한 행이 있으면 바로 반환한다. 없으면 행 거리의 절댓값으로 정렬한 가장 가까운 두 점 $a,b$로 다음을 계산한다.
\[
 \hat c_r=\clip\left(c_a+\frac{(r-a)(c_b-c_a)}{b-a},0,83\right).
\]
두 점이 없으면 None이다. 두 점이 모두 같은 쪽에 있으면 외삽이 된다. 현재 선택 모드에서는 42행 누락을 보완하는 데 사용한다.

\section{장애물 검출과 회피 방향}
\subsection{밝은 연결 영역 찾기}
도로 중심이 3개 미만이면 장애물을 찾지 않는다. 밝기 $I\ge0.54$인 픽셀 중 $22\le y<62$만 조사한다. 배열 전체 크기의 방문 표시를 만들고, 위에서 아래·왼쪽에서 오른쪽 순서로 픽셀을 탐색한다. 아직 방문하지 않은 밝은 픽셀마다 스택을 사용한 8방향 연결 영역 탐색을 수행한다.

각 영역의 면적, 최소·최대 좌표, 좌표 합을 계산한다. 다음 조건을 모두 만족해야 후보가 된다.
\begin{itemize}
\item 면적: 4 이상 80 이하 픽셀
\item 바운딩 박스 너비: 2 이상 9 이하 픽셀
\item 바운딩 박스 높이: 2 이상 10 이하 픽셀
\item 영역 중심 $x$와 그 높이의 도로 중심 사이: 12픽셀 이하
\end{itemize}
영역 중심은 면적에 포함된 픽셀 좌표의 평균이다. 최종 선택은 중심 $y$가 가장 큰 후보다. 반환값은 $(y_o,x_o,c_o)$ 또는 None이다. 실제 거리·장애물 ID·종류를 출력하지 않는다. 가장 아래에 보이는 후보를 가장 가까운 것으로 취급한다.

\subsection{회피 방향 상태의 모든 조건}
장애물의 도로 상대 위치를 $d=x_o-c_o$라 하자. $d<0$이면 회피 부호 $q=+1$, 그 외에는 $q=-1$을 후보로 삼는다. 실제 코드의 조향 부호로 양수 쪽 또는 음수 쪽에 회피 항을 더한다.

\begin{enumerate}
\item 저장된 회피 부호가 0이면 새 후보를 채택하고 조우 횟수를 1 늘린다.
\item 부호가 이미 있고 $y_o<52$이면 방향을 다시 판단한다. $|d|>1$이면 후보 부호를 바로 채택한다.
\item $|d|\le1$이고 이전 상대 위치가 있으면 이동량 $\Delta d$를 본다. 기존 $q>0$, $d>-3$, $\Delta d>2$이면 $q=-1$로 바꾼다. 기존 $q<0$, $d<3$, $\Delta d<-2$이면 $q=+1$로 바꾼다.
\item $y_o\ge52$에서는 이미 정한 부호를 재평가하지 않는다. 상대 위치 기록은 갱신한다.
\item 검출되면 누락 횟수를 0으로 초기화하고 검출 횟수를 늘린다.
\item 검출이 사라지면 첫 누락에는 부호를 유지한다. 두 번째 연속 누락에서 부호·이전 상대 위치를 지운다. 누락된 프레임에는 회피 조향 항 자체를 더하지 않는다.
\end{enumerate}

장애물 긴급도와 회피 조향 항은
\[
 u=\clip((y_o-22)/18,0,1),\qquad s_o=q\cdot0.34\cdot u
\]
이다. $y_o\ge40$이면 긴급도는 이미 1이다. 현재 회피 배율은 1.0이다. 이 구조는 장애물 ID를 연결하는 추적기가 아니며, 다음 프레임에서 선택된 다른 물체에도 같은 상태가 사용될 수 있다.

\section{현재 속도 추정}
\code{SPEED_ROI=(77,83,10,13)}을 슬라이스로 사용하므로 실제 영역은 행 77--82, 열 10--12다. 밝기 합을 $M$이라 할 때 한 프레임의 추정 속도는
\[
 v_{\mathrm{frame}}=\clip((M-0.27)/0.085,0,80).
\]
네 장의 관측에서는 마지막 두 장 각각을 디코딩한 뒤 평균한다. 단일 화면이면 그 화면의 값 하나를 사용한다. 이 평균은 흔들림을 줄일 수 있지만 최근 변화에 지연을 더한다.

80은 이 \textbf{속도 추정 함수의 클리핑 상한}이다. 게임 엔진의 최고속도와 같은 개념이 아니다. 이 제출 코드에는 상한을 넘는 HUD를 별도로 해석하는 함수가 없다.

\section{하위 도로 추종기: VisionCorridorAgent}
\subsection{생성 파라미터와 유효성 검사}
기본값과 이 제출 경로가 실제로 전달하는 값은 다르다.
\par\noindent\begin{tabularx}{\linewidth}{@{}lrrX@{}}\toprule
항목 & 기본값 & 실제 값 & 허용 범위 \\\midrule
cruise speed & 62 & 55 & 유한값 36--80 \\
curve penalty & 2 & 3 & 유한값 0--4 \\
max gas & 0.12 & 0.5 & 유한값 $(0,1]$ \\
obstacle steer scale & 1 & 1 & 유한값 0--2 \\
far obstacle speed & 47 & 43 & 유한값 20--80 \\
near obstacle speed & 40 & 36 & 유한값 20--80 \\\bottomrule
\end{tabularx}
범위를 벗어나면 ValueError가 발생한다. 클래스의 도로·속도 밝기 상수와 같은 이름의 모듈 상수가 있지만, 실제 인식 래퍼는 \code{pixel_features} 함수를 호출한다. 클래스 상수만 바꿔도 인식 함수의 상수가 함께 바뀌는 구조는 아니다.

\subsection{기본 조향과 하위 목표 속도}
없는 행은 42를 대신 사용한다. $m=c_{42}$, $n=c_{54}$라 하면
\[
 s_r=0.022(m-42)+0.018(m-n).
\]
도로 추종 항에 검출된 장애물 회피 항을 더하고 $[-0.7,0.7]$로 자른다. 즉 $s_{\mathrm{base}}=\clip(s_r+s_o,-0.7,0.7)$이다.

하위 \code{road_sweep}는 $r\in\{30,34,38,42\}$에 대해 $|c_r-c_{r+24}|$의 최댓값을 계산한다. 없는 행은 42다. 탐색 행에 58, 62, 66은 없으므로 해당 쌍의 두 번째 값은 42다. 이는 상위의 전체 중심 최댓값-최솟값과 다른 함수다.

하위 목표는 $\clip(55-3\,\mathrm{road\_sweep},36,55)$이다. 장애물 높이 $y_o<44$면 43, 그 외에는 36으로 상한을 낮춘다. 장애물 첫 누락에도 43 상한을 둔다. 이전 목표가 있으면 새 목표에 $0.65T_{\mathrm{old}}+0.35T_{\mathrm{new}}$ 평활화를 적용한다.

\subsection{하위 페달과 실제 효과}
하위 함수는 $v>T+1$일 때 가속 0, 제동 $\clip(0.012(v-T),0.04,0.28)$을 반환한다. 그렇지 않으면 $v<T-8$에서 gas=0.5, $v<T-3$에서 gas=$0.5\cdot2/3$, 나머지에서 gas=$0.5\cdot5/12$이며 brake=0이다.

\textbf{하지만 이 하위 가속·제동은 현재 제출 경로에서 매번 상위 함수로 덮어써진다.} 최종 가속이 최대 0.5라고 해석하면 틀린다. 하위 조향, 장애물 상태와 진단 카운터는 사용되지만, 하위 목표의 평활화·43/36 제한·페달 분기는 최종 페달을 직접 결정하지 않는다.

\section{조향 보정과 고정 속도 계층}
\subsection{실제로 선택된 damping 모드}
상위 \code{FixedHighSpeedAgent}는 하위 act를 호출하고 그 조향에 보정 $c$를 더해 다시 자른다. 11번째 act부터 현재 30·42·54행이 모두 있고 이전 화면에도 42행이 있으면
\[
 c=0.025\bigl(c_{42}^{\mathrm{current}}-c_{42}^{\mathrm{previous}}\bigr).
\]
조건을 만족하지 못하면 $c=0$이다. 모드 이름은 damping이지만 실제 계산은 화면 중심의 프레임 간 변화에 대한 가산 항이다. 차량의 실제 요레이트나 타이어 슬립을 측정하는 제어는 아니다.

중요한 순서는 $\clip(\clip(s_r+s_o,-0.7,0.7)+c,-0.7,0.7)$이다. 두 번 클리핑하므로 모든 항을 한 번에 더하고 자른 결과와 다를 수 있다.

\subsection{첫 10회와 고정 페달}
steps는 reset에서 0이다. 보정 조건은 증가 전 steps가 10 이상이고, 증가 연산은 고정 계층 act의 끝에 있다. 따라서 첫 10번 호출에는 시간 보정이 없고, 11번째부터 적용된다. 기본 도로 조향과 장애물 회피는 첫 호출부터 동작한다.

고정 계층은 매번 다음 페달을 만든다.
\[
 g_{60}=\clip(0.12+0.04(60-v),0,0.6),\quad
 b_{60}=\clip(0.02(v-62),0,0.15).
\]
추정 속도가 0이면 gas=0.6, brake=0이다. 현재 제출 코드에는 gas=1로 시작하는 별도 출발 모드가 없다.

\section{급감속과 충돌 가능성 처리}
선택된 회복 모드는 \code{impact_clear}다. 이전 최대 4개 추정 속도를 저장한다. 다음 조건을 모두 만족하면 충돌 가능성을 감지한다.
\begin{itemize}
\item 증가 후 steps가 10 초과
\item 현재 impact 카운터가 0
\item 속도 이력이 비어 있지 않음
\item 현재 추정 속도 $v<50$
\item 이전 4개 이력의 최댓값에서 현재 속도를 뺀 값이 10 초과
\end{itemize}
감지되면 impact 카운터를 12로 설정한다. 도로 42·54행이 모두 있고 카운터가 양수이면 장애물 회피 항을 빼고 $\clip(s_r+c,-0.7,0.7)$로 조향을 다시 계산한다. 카운터는 해당 호출 끝에서 1 감소하므로, 감지한 호출을 포함해 최대 12회 조건부 회피 해제를 적용한다. 도로가 누락된 호출에도 카운터 시간은 소비된다.

이는 실제 충돌 플래그를 읽지 않는 속도 기반 추정이다. 이 계층 자체는 가속·제동을 바꾸지 않는다. 실제로 멈췄는지 판단한 뒤 후진하거나 차량을 재배치하는 기능도 없다.

\section{최상위 조정: preview\_row\_repair}
\subsection{정상 제동을 충돌로 오인하지 않기}
위의 충돌 감지가 \textbf{이번 호출에서 새로 발생}했고 이전 최대 4개 최종 brake 중 하나라도 0.01보다 크면 감지를 취소한다. 조향을 회복 계층 적용 전의 \code{damping_steer}로 되돌리고 impact 카운터를 0으로 만든다. 현재 호출의 brake는 나중에 계산되므로 이 판단에는 아직 들어가지 않는다.

이 취소는 기존 impact 상태 전체를 매번 재검토하는 규칙이 아니다. 새 trigger가 생긴 순간에만 작동한다. 정당한 제동과 충돌이 가까운 시간에 함께 일어나면 실제 충돌도 가려질 가능성이 있다.

\subsection{42행이 사라진 경우}
steps $>10$, 42행 없음, 54행 있음, 알려진 중심 최소 2개일 때 \code{observed_center}로 42행을 복원한다. 이 복원값과 54행, 현재 검출된 장애물 회피 항, 시간 보정으로 조향을 한 번에 다시 계산하고 $[-0.7,0.7]$로 자른다.

복원된 중심은 원래 centers 사전에 삽입되지 않는다. 따라서 바로 뒤의 속도 계산은 원래 검출된 점만 사용한다. 이 모드에서는 54행까지 사라지면 같은 복원 분기가 작동하지 않는다. 모든 도로를 잃었을 때 장기간 방향을 기억하는 기능은 다른 선택 모드에만 있다.

\subsection{최종 목표 속도와 페달}
첫 10번 호출은 앞 계층의 60 목표 페달을 유지한다. 이후에는 검출된 도로 중심의 범위
\[
 S=\begin{cases}\max c_y-\min c_y,&\text{중심이 2개 이상},\\28,&\text{그 외}\end{cases}
\]
로 $T=\max(38,60-0.8S)$를 계산한다. 장애물이 현재 검출되면 $T\leftarrow\min(T,44)$이다. 중심이 거의 없으면 $60-0.8\cdot28=37.6$에서 바닥 제한으로 38이 된다. 장애물은 44보다 더 낮은 도로 목표를 올려 주지는 않는다.

$T<60$이면 최종 페달을
\[
 g=\clip(0.12+0.04(T-v),0,0.6),\quad
 b=\clip(0.02(v-T-2),0,0.28)
\]
로 덮어쓴다. $T=60$이면 기존의 brake 최대 0.15를 유지한다. 목표가 60보다 아주 조금 낮아지는 순간에도 제동 상한이 0.28로 바뀐다는 불연속이 있다.

가속·제동을 각각 독립 계산하므로 $T+2<v<T+3$에서 둘 다 양수가 될 수 있다. 목표와 속도가 같아도 gas=0.12다. 이는 목표 속도를 유지하려는 가산항이지 물리적으로 교정된 유지 토크는 아니다.

\subsection{계산 예시}
\par\noindent\begin{tabularx}{\linewidth}{@{}Xrrrr@{}}\toprule
조건 & 목표 & 속도 & gas & brake \\\midrule
직선, 충분히 느림 & 60 & 40 & 0.60 & 0 \\
직선, 목표 도달 & 60 & 60 & 0.12 & 0 \\
직선, 초과 속도 & 60 & 70 & 0 & 0.15 \\
장애물 제한 & 44 & 60 & 0 & 0.28 \\
작은 초과 속도 & 44 & 46.5 & 0.02 & 0.01 \\
큰 도로 중심 범위 & 38 & 45 & 0 & 0.10 \\\bottomrule
\end{tabularx}
표는 코드 식을 대입한 값이다. 이 명령이 만드는 실제 가속도나 정지거리를 측정한 표는 아니다.

\section{한 번의 act에서 실제로 실행되는 순서}
\begin{enumerate}
\item 최상위가 호출 전 impact 카운터를 보관한다. 선택 모드의 최종 판단에는 이 값이 직접 사용되지 않는다.
\item 고정 계층이 현재 프레임을 검증하고 도로 중심을 찾는다.
\item 내부 base가 다시 프레임·도로를 읽고 장애물을 검출한다. 회피 부호와 자체 목표·진단 상태를 갱신하고 행동을 반환한다.
\item 고정 계층이 시간 보정 조향과 60 목표 페달을 적용하고 steps를 증가시킨다.
\item 회복 계층이 급감속을 감지하고 필요하면 회피 조향을 해제한다. 속도 이력과 impact 카운터를 갱신한다.
\item 최상위가 최근 제동에 의한 감지 취소를 적용한다.
\item 필요한 경우 42행을 복원하고 조향을 다시 계산한다.
\item 도로 범위와 현재 장애물로 최종 목표 속도를 정해 페달을 덮어쓴다.
\item 마지막 조향, 최대 4개 최종 제동 이력, 진단 사전을 갱신하고 행동을 반환한다.
\end{enumerate}
도로 검출과 속도 디코딩 일부가 중복 호출된다. 이 코드에는 이를 한 번만 계산해 공유하는 별도의 캐시 계층이 없다.

\section{상태 변수와 reset}
reset은 전달된 화면으로 위치를 학습하지 않고 상태를 초기화한다. 하위 base는 observation을 버린다. 상속 체인의 reset이 차례로 실행되어 다음 상태를 지운다.
\begin{longtable}{@{}p{69mm}p{91mm}@{}}\toprule
상태 & 초기값과 역할 \\\midrule\endhead
\code{_obstacle_side}, \code{_obstacle_missing} & 0, 0. 회피 방향과 연속 검출 누락 수 \\
\code{_last_obstacle_side_offset} & None. 이전 장애물의 도로 상대 위치 \\
\code{_target_speed} & None. 하위 목표 평활화 기록 \\
\code{_last_obstacle} & None. 현재 검출값; 각 act에서도 먼저 None으로 지움 \\
\code{_obstacle_encounters}, \code{_obstacle_detection_frames} & 0. 조우·검출 집계 \\
\code{_speed_samples}, \code{_speed_sum}, \code{_speed_max} & 0. 하위 속도 추정 집계 \\
\code{_gas_frames}, \code{_brake_frames} & 0. 하위 페달이 양수였던 횟수 \\
\code{_last_speed_estimate}, \code{_last_target_speed}, \code{_last_gas}, \code{_last_brake} & 0. 하위 마지막 판단 기록 \\
\code{steps}, \code{last} & 0, 빈 사전. 호출 수와 최상위에서 반환할 진단 \\
\code{remembered_road}, \code{latched_turn}, \code{recovery_left}, \code{recovery_armed} & 0, 0, 0, True. 기억 복귀 모드용; 현재 선택 모드에서는 복귀 명령을 켜지 않음 \\
\code{centered_streak} & 0. 7개 중심이 있고 42·54행 모두 화면 중심에서 8 이내이면 증가. 3회면 회복 카운터 초기화·재무장 \\
\code{impact_left}, \code{speed_history} & 0, 빈 리스트. 급감속 후 회피 해제 기간과 최근 4개 속도 \\
\code{last_command}, \code{brake_history} & 0, 빈 리스트. 마지막 최종 조향과 최근 4개 최종 brake \\
\code{recovering}, \code{stable}, \code{visible_steer}, \code{lost_rows} & False, 0, 0, 0. 다른 조정 모드용 상태 \\\bottomrule
\end{longtable}
생성자 안의 \code{self.reset()}은 파이썬의 동적 메서드 호출에 의해 최상위 reset까지 실행된다. 각 reset이 super를 호출하므로 하위 상태도 초기화된다.

\section{진단 정보의 의미와 주의점}
Agent에서 노출하는 진단은 \code{dict(self.last)}다. 하위 base의 진단 함수는 별도로 존재하며 이 반환값에 자동 합쳐지지 않는다.
\begin{itemize}
\item 고정 계층: target speed, pixel speed, inherited steer, correction, changed, available, road centers, final gas/brake를 기록한다.
\item 회복 계층: recovery mode, damping steer, recovery changed, recovery steps remaining, centered streak, impact proxy trigger, impact steps remaining을 추가한다.
\item 최상위: coordination mode, 최종 target speed, braking proxy veto, geometry repaired, completion changed, 최종 gas/brake를 덮어쓰거나 추가한다.
\end{itemize}
\code{changed}는 고정 계층 보정 전후 차이, \code{recovery_changed}는 회복 계층 조향 변경, \code{completion_changed}는 최상위와 회복 계층 반환 행동의 차이다. 어느 것도 전체 제출 전략과 별도 대조군 사이의 차이를 뜻하지 않는다.

제동 오인으로 impact를 취소하면 실제 \code{impact_left}는 0이 되지만 이미 기록된 \code{impact_steps_remaining}은 다시 고쳐지지 않는다. \code{braking_proxy_veto}와 실제 상태를 함께 해석해야 한다. 하위 gas/brake 집계는 상위가 덮어쓰기 전 수치이므로 최종 페달 사용률이 아니다.

base의 \code{diagnostics()}는 장애물 조우/검출, 속도 평균/최대, gas/brake 프레임 수를 반환한다. base의 \code{last_step_diagnostics()}는 하위 target, speed, gas, brake와 장애물 좌표/도로 중심/회피 부호를 반환한다. 값이 없으면 일부 항목은 None이다.

\section{포함되어 있지만 선택되지 않은 모드}
\subsection{FixedHighSpeedAgent의 다른 조향 모드}
\begin{itemize}
\item control: 시간 또는 앞보기 보정 없이 하위 조향을 쓴다.
\item preview: 보정 $0.018(c_{30}-c_{42})$를 더한다.
\item curvature: 보정 $0.018(c_{30}-2c_{42}+c_{54})$를 더한다.
\item pursuit: 원하는 조향을 $0.65\,[48(c_{30}-42)/((c_{30}-42)^2+24^2)]$로 정하고, 기본 road 항을 뺀 값을 보정으로 더한다. 화면 좌표 식이며 미터 단위 카메라 교정 주장은 없다.
\end{itemize}
이 모드들도 보정은 첫 10회 이후, 30·42·54행이 있을 때만 작동한다. 잘못된 모드 문자열은 ValueError다. 현재 인스턴스는 damping으로 고정된다.

\subsection{memory\_reacquire}
42행이 사라지고 회복이 재무장된 상태이며 기억한 도로 조향의 절댓값이 0.01보다 크면, 그 방향으로 크기 0.35--0.7의 조향을 고정한다. 카운터는 80이다. 7개 중심이 보이고 42·54행이 중심에서 8 이내인 상태가 3회 이어지면 종료하고 다시 무장한다. 현재 선택된 impact\_clear 모드에서는 이 분기를 실행하지 않는다.

\subsection{FastCompletionCoordination의 나머지 모드}
모든 모드는 하위 impact\_clear와 damping을 포함한다. 이후 조정만 다르다.
\begin{longtable}{@{}p{49mm}p{111mm}@{}}\toprule
모드 & 추가 동작 \\\midrule\endhead
control & 최상위 추가 제어 없이 60 목표를 유지 \\
preview\_budget & 도로 중심 범위 및 장애물로 목표 38--60, 장애물 최대 44 \\
preview\_action\_aware & preview budget + 최근 제동에 의한 급감속 감지 취소 \\
actuator\_rate & 이전 최종 명령에서 조향 변화량을 호출당 ±0.24로 제한 \\
impact\_traction & 호출 전 impact가 있거나 이번에 trigger가 발생하면 gas를 최대 0.18로 제한 \\
recovery\_budget & 54행 중심 오차가 6 초과 또는 42행 누락이면 회복 상태. 목표 40. 7개 중심과 오차 4 미만이 3회 연속이면 해제 \\
recovery\_geometry\_aware & 제동 오인 취소 + 양쪽 중심 복원/조향 기억 + recovery budget \\
preview\_row\_repair & 실제 선택. preview action aware + 조건부 42행 복원 \\\bottomrule
\end{longtable}
recovery geometry aware는 42 또는 54행이 없으면 두 행을 각각 보간·외삽한다. 둘 다 구해지면 도로항, 회피항, 이전 42행과의 변화 $0.025(m-pm)$을 합친다. impact 중이며 취소되지 않았다면 회피항을 0으로 한다. 재계산한 조향을 기억한다. 두 점을 구할 수 없으면 최대 12회 기억 조향을 유지한다. 두 행이 정상적으로 보이면 현재 조향을 기억하고 누락 수를 초기화한다. 회복 예산의 54행 누락 기본값은 0이어서 오차가 42가 된다.

이 모드 이름들이 파일에 존재한다는 사실은 현재 ZIP에서 모두 실행된다는 뜻이 아니다. 선택은 agent.py와 생성자 체인이 결정한다.

\section{사용되지 않는 특징 벡터 함수도 포함됨}
\subsection{7차원 시각 특징}
\code{extract_visual_features}는 다음 값을 float32 배열로 반환한다. 현재 Agent의 act는 이 함수를 호출하지 않는다.
\begin{enumerate}
\item 평균 HUD 속도 / 80
\item $(c_{54}-42)/42$, $[-1,1]$로 클리핑
\item $(c_{42}-42)/42$, $[-1,1]$로 클리핑
\item 하위 road sweep / 42, $[0,1]$로 클리핑
\item 장애물 존재: 1 또는 0
\item 장애물 상대 위치 $(x_o-c_o)/12$, $[-1,1]$로 클리핑
\item 긴급도 $\clip((y_o-22)/18,0,1)$
\end{enumerate}
장애물이 없으면 마지막 세 값은 모두 0이다. FEATURE\_NAMES는 이 순서의 문자열 이름을 보관하고 VISUAL\_FEATURE\_SIZE는 7이다.

\subsection{1차원 시간 특징}
\code{extract_temporal_features}는 $(4,84,84)$ 모양을 요구한다. 현재 장애물이 없거나 현재 정규화 상대 위치의 절댓값이 0.2 이상이면 0이다. 이전 화면에서 장애물을 찾지 못해도 0이다. 이전 상대 위치를 /12 후 $[-1,1]$로 자르고, 그 절댓값이 0.3 미만이면 0, 그 외에는 이전 상대 위치 하나를 반환한다. 이름은 previous\_obstacle\_lateral\_offset이고 크기는 1이다. 이 함수도 현재 주행에서 사용되지 않는다.

\subsection{배치 처리와 export 목록}
두 batch 함수는 $(B,4,84,84)$ 모양과 전체 값의 유한성·$[0,1]$ 범위를 검사하고 각 관측의 결과를 \code{np.stack}한다. 빈 배치를 별도로 처리하지 않으므로 빈 입력에 대한 정상 반환을 보장하지 않는다. 시간 특징의 단일 함수는 모양 검사만 명시적으로 수행한다.

\code{pixel_features.__all__}은 특징 이름/크기와 인식·특징 추출 함수들을 공개 이름으로 열거한다. \code{corridor_agent.__all__}은 VisionCorridorAgent 하나다. 이 목록은 실행 순서나 활성 모드를 선택하지 않는다. 파일 첫 설명의 ``learned policy''라는 문장은 보조 특징 모듈의 설명이며, 이 제출 ZIP에 학습 정책이 존재한다는 증거가 아니다.

\section{예외, 미구현 기능, 실제 전략의 경계}
\begin{itemize}
\item base 단독의 \code{_frame}은 ValueError를 잡아 None을 반환하고 base act는 영벡터를 돌려준다. 그러나 현재 상위 act가 먼저 current\_frame을 호출하므로 최종 Agent 전체가 항상 영벡터로 복구한다고 말할 수 없다.
\item 도로 검출 누락은 해당 행 생략, 기본 중심 42, 조건부 42행 복원으로 처리한다. 모든 상황을 위한 탐색 주행은 없다.
\item 장애물이 없거나 검출 조건에서 탈락하면 None이다. 별도 신뢰도나 여러 장애물의 공동 경로 계획은 없다.
\item out-in-out 경로 생성, 전체 맵 기억, 곡률의 물리 단위 추정, 제동거리 기반 속도 계획, 모델 예측 제어는 이 ZIP에 없다.
\item 실제 차량 속도·충돌·좌표를 환경에서 직접 읽지 않는다. 입력은 화면이고 명령 이력 및 내부 상태만 기억한다.
\item 후진·재시작·순간이동·완주 여부 판정은 이 Agent의 기능이 아니다. 완주와 차량 물리는 외부 환경이 처리한다.
\item 이 패키지는 실행 중 학습하거나 가중치를 갱신하지 않는다. 같은 관측과 상태에서는 동일한 계산식을 따른다.
\end{itemize}

\section{최종 요약}
최종 행동은 단일 식 하나가 아니라 하위 도로 추종, 시간 보정, 급감속 회피 해제, 제동 오인 취소, 누락 도로 복원, 최종 속도 조정이 순서대로 만든 결과다. 하위에서 계산한 행동 중 일부가 덮어써지므로, 코드에 존재하는 숫자와 실제 최종 조작에 적용되는 숫자를 구별해야 한다.

현재 전략을 한 문장으로 표현하면 다음과 같다: \textbf{픽셀 도로 중심을 따라 조향하고 가장 가까워 보이는 장애물의 반대쪽으로 회피 항을 더하며, 보이는 도로 중심의 좌우 범위로 속도 목표를 낮추고, 급감속·인식 누락 때 조향을 복구하는 규칙 기반 제어기다.}

\clearpage\appendix
\section{패키지 식별 정보}

ZIP SHA-256:\par
{\small\ttfamily 7bc38a3bdeab7cfdf437cfbf5c06f15a81cfa18fad5317748ab13a259b8ea835}

ZIP 크기: 9928 bytes. 파일 수: 7.
\begin{longtable}{@{}p{75mm}rr@{}}\toprule 파일 & bytes & 줄 수 \\\midrule\endhead
\code{agent.py} & 388 & 6 \\
\code{haic_agent/__init__.py} & 32 & 1 \\
\code{haic_agent/fixed_high_speed_runtime.py} & 3105 & 70 \\
\code{haic_agent/corridor_agent.py} & 10439 & 258 \\
\code{haic_agent/pixel_features.py} & 9729 & 247 \\
\code{haic_agent/fixed_high_speed_recovery_v2.py} & 2906 & 61 \\
\code{haic_agent/fast_completion_coordination.py} & 5285 & 98 \\
\bottomrule\end{longtable}

\clearpage
\section{agent.py 전체 소스}
{\footnotesize SHA-256:\par\ttfamily a58359bb8469c60cff05e78d2749ff5e85287fabfc1756e064016f13aa69e579}
\begin{lstlisting}
from haic_agent.fast_completion_coordination import FastCompletionCoordination
class Agent:
    def __init__(self): self.driver=FastCompletionCoordination("preview_row_repair")
    def reset(self,observation): self.driver.reset(observation)
    def act(self,observation): return self.driver.act(observation)
    def last_step_diagnostics(self): return self.driver.last_step_diagnostics()

\end{lstlisting}

\clearpage
\section{haic\_agent/\_\_init\_\_.py 전체 소스}
{\footnotesize SHA-256:\par\ttfamily a5e2b46bb2f58049e8c627b3f3bd75f6b80ddbd6e8e9270b1b6dca19edee01d8}
\begin{lstlisting}
"""Pixel controller runtime."""

\end{lstlisting}

\clearpage
\section{haic\_agent/fixed\_high\_speed\_runtime.py 전체 소스}
{\footnotesize SHA-256:\par\ttfamily ab465ff850fb489ff5facde49b8979a2e1e164885015703edf33b81c531e8e88}
\begin{lstlisting}
"""Pixel-only fixed-speed steering mechanisms for a TRAIN feasibility study."""
import numpy as np

from haic_agent.corridor_agent import VisionCorridorAgent
from haic_agent.pixel_features import current_frame, road_centers, estimate_observation_speed


ARMS = ("control", "preview", "curvature", "damping", "pursuit")
TARGET_SPEED = 60.0
PREFIX_DECISIONS = 10


def fixed_pedals(speed):
    return (float(np.clip(0.12 + 0.04 * (TARGET_SPEED - speed), 0, 0.6)),
            float(np.clip(0.02 * (speed - TARGET_SPEED - 2), 0, 0.15)))


class FixedHighSpeedAgent:
    def __init__(self, mechanism="control"):
        if mechanism not in ARMS:
            raise ValueError("unknown steering mechanism")
        self.mechanism = mechanism
        self.base = VisionCorridorAgent(cruise_speed=55, curve_speed_penalty=3, max_gas=0.5,
                                        obstacle_far_speed=43, obstacle_near_speed=36)
        self.reset()

    def reset(self, observation=None):
        self.base.reset(observation)
        self.steps = 0
        self.last = {}

    def act(self, observation):
        frame = current_frame(observation)
        centers = road_centers(frame)
        inherited = self.base.act(observation)
        near, middle = centers.get(54, 42.0), centers.get(42, 42.0)
        baseline_road = 0.022 * (middle - 42) + 0.018 * (middle - near)
        correction = 0.0
        available = all(row in centers for row in (30, 42, 54))
        if self.steps >= PREFIX_DECISIONS and available:
            far = centers[30]
            if self.mechanism == "preview":
                correction = 0.018 * (far - middle)
            elif self.mechanism == "curvature":
                correction = 0.018 * (far - 2 * middle + near)
            elif self.mechanism == "damping":
                pixels = np.asarray(observation)
                previous = road_centers(pixels[-2]) if pixels.ndim == 3 else {}
                if 42 in previous:
                    correction = 0.025 * (middle - previous[42])
                else:
                    available = False
            elif self.mechanism == "pursuit":
                # Image-space bearing; no claim of metric camera calibration.
                desired = 0.65 * (2 * 24 * (far - 42) / ((far - 42) ** 2 + 24 ** 2))
                correction = desired - baseline_road
        steer = float(np.clip(float(inherited[0]) + correction, -self.base.MAX_STEER,
                              self.base.MAX_STEER))
        speed = estimate_observation_speed(observation, frame)
        gas, brake = fixed_pedals(speed)
        self.steps += 1
        self.last = dict(target_speed=TARGET_SPEED, pixel_speed=float(speed),
                         inherited_steer=float(inherited[0]), correction=correction,
                         changed=abs(steer - float(inherited[0])) > 1e-6,
                         available=available, road_centers=centers,
                         final_gas=gas, final_brake=brake)
        return np.asarray([steer, gas, brake], dtype=np.float32)

    def last_step_diagnostics(self):
        return dict(self.last)

\end{lstlisting}

\clearpage
\section{haic\_agent/corridor\_agent.py 전체 소스}
{\footnotesize SHA-256:\par\ttfamily 00667a41e6ed07c715784669745dcf4646307214474daf2e87b73ee48dc835cd}
\begin{lstlisting}
"""Pixel-only racing controller using road geometry, speed HUD, and obstacle cues."""

from __future__ import annotations

import numpy as np

from haic_agent.pixel_features import (
    center_at as pixel_center_at,
    current_frame as pixel_current_frame,
    estimate_observation_speed as pixel_estimate_observation_speed,
    estimate_speed as pixel_estimate_speed,
    nearest_bright_object as pixel_nearest_bright_object,
    road_centers as pixel_road_centers,
    road_sweep as pixel_road_sweep,
)


class VisionCorridorAgent:
    """Follow the road corridor and regulate speed from rendered pixels only."""

    MAX_STEER = 0.7
    ROAD_LOW = 0.24
    ROAD_HIGH = 0.52
    OBSTACLE_LOW = 0.54
    OBSTACLE_STEER = 0.34
    MAX_TEACHER_GAS = 0.12
    MAX_TEACHER_BRAKE = 0.28
    SPEED_ROI = (77, 83, 10, 13)
    SPEED_BASELINE = 0.27
    SPEED_PER_UNIT = 0.085

    def __init__(
        self,
        *,
        cruise_speed: float = 62.0,
        curve_speed_penalty: float = 2.0,
        max_gas: float = MAX_TEACHER_GAS,
        obstacle_steer_scale: float = 1.0,
        obstacle_far_speed: float = 47.0,
        obstacle_near_speed: float = 40.0,
    ) -> None:
        if not np.isfinite(cruise_speed) or not 36.0 <= cruise_speed <= 80.0:
            raise ValueError("cruise_speed must be finite and between 36 and 80")
        if not np.isfinite(curve_speed_penalty) or not 0.0 <= curve_speed_penalty <= 4.0:
            raise ValueError("curve_speed_penalty must be finite and between 0 and 4")
        if not np.isfinite(max_gas) or not 0.0 < max_gas <= 1.0:
            raise ValueError("max_gas must be finite and in (0, 1]")
        if not np.isfinite(obstacle_steer_scale) or not 0.0 <= obstacle_steer_scale <= 2.0:
            raise ValueError("obstacle_steer_scale must be finite and between 0 and 2")
        for name, value in (
            ("obstacle_far_speed", obstacle_far_speed),
            ("obstacle_near_speed", obstacle_near_speed),
        ):
            if not np.isfinite(value) or not 20.0 <= value <= 80.0:
                raise ValueError(f"{name} must be finite and between 20 and 80")
        self.cruise_speed = float(cruise_speed)
        self.curve_speed_penalty = float(curve_speed_penalty)
        self.max_gas = float(max_gas)
        self.obstacle_steer_scale = float(obstacle_steer_scale)
        self.obstacle_far_speed = float(obstacle_far_speed)
        self.obstacle_near_speed = float(obstacle_near_speed)
        self._obstacle_side = 0.0
        self._obstacle_missing = 0
        self._last_obstacle_side_offset: float | None = None
        self._target_speed: float | None = None
        self._obstacle_encounters = 0
        self._obstacle_detection_frames = 0
        self._speed_samples = 0
        self._speed_sum = 0.0
        self._speed_max = 0.0
        self._gas_frames = 0
        self._brake_frames = 0
        self._last_obstacle: tuple[float, float, float] | None = None
        self._last_speed_estimate = 0.0
        self._last_target_speed = 0.0
        self._last_gas = 0.0
        self._last_brake = 0.0

    def reset(self, observation=None) -> None:
        """Clear route-dependent state between episodes."""
        del observation
        self._obstacle_side = 0.0
        self._obstacle_missing = 0
        self._last_obstacle_side_offset = None
        self._target_speed = None
        self._obstacle_encounters = 0
        self._obstacle_detection_frames = 0
        self._speed_samples = 0
        self._speed_sum = 0.0
        self._speed_max = 0.0
        self._gas_frames = 0
        self._brake_frames = 0
        self._last_obstacle = None
        self._last_speed_estimate = 0.0
        self._last_target_speed = 0.0
        self._last_gas = 0.0
        self._last_brake = 0.0

    @staticmethod
    def _frame(observation) -> np.ndarray | None:
        try:
            return pixel_current_frame(observation)
        except ValueError:
            return None

    def _road_centers(self, frame: np.ndarray) -> dict[int, float]:
        return pixel_road_centers(frame)

    @staticmethod
    def _center_at(row: float, centers: dict[int, float]) -> float:
        return pixel_center_at(row, centers)

    def _nearest_bright_object(
        self, frame: np.ndarray, centers: dict[int, float]
    ) -> tuple[float, float, float] | None:
        """Find the nearest small bright component inside the tracked road band."""
        return pixel_nearest_bright_object(frame, centers)

    @classmethod
    def _estimate_speed(cls, frame: np.ndarray) -> float:
        """Decode speed from the lower HUD bar at its calibrated pixel scale."""
        return pixel_estimate_speed(frame)

    def _observation_speed(self, observation, current_frame: np.ndarray) -> float:
        return pixel_estimate_observation_speed(observation, current_frame)

    @staticmethod
    def _road_sweep(centers: dict[int, float]) -> float:
        """Estimate upcoming bend strength from the pixel centerline spread."""
        return pixel_road_sweep(centers)

    def _pedals(self, speed: float, target_speed: float) -> tuple[float, float]:
        if speed > target_speed + 1.0:
            brake = float(np.clip((speed - target_speed) * 0.012, 0.04, self.MAX_TEACHER_BRAKE))
            return 0.0, brake
        if speed < target_speed - 8.0:
            return self.max_gas, 0.0
        if speed < target_speed - 3.0:
            return self.max_gas * (2.0 / 3.0), 0.0
        return self.max_gas * (5.0 / 12.0), 0.0

    def act(self, observation) -> np.ndarray:
        frame = self._frame(observation)
        if frame is None:
            return np.zeros(3, dtype=np.float32)

        centers = self._road_centers(frame)
        self._last_obstacle = None
        far = centers.get(42, 42.0)
        near = centers.get(54, 42.0)
        steering = 0.022 * (far - 42.0) + 0.018 * (far - near)
        target_speed = float(
            np.clip(
                self.cruise_speed - self.curve_speed_penalty * self._road_sweep(centers),
                36.0,
                self.cruise_speed,
            )
        )

        obstacle = self._nearest_bright_object(frame, centers)
        if obstacle is not None:
            self._last_obstacle = obstacle
            obstacle_y, obstacle_x, road_center = obstacle
            # The image side of an obstacle can shift with perspective and bends.
            # Re-evaluate to the near field, ignoring sub-pixel side noise.
            if self._obstacle_side == 0.0:
                self._obstacle_encounters += 1
            side_offset = obstacle_x - road_center
            candidate_side = 1.0 if side_offset < 0.0 else -1.0
            if self._obstacle_side == 0.0:
                self._obstacle_side = candidate_side
            elif obstacle_y < 52.0:
                if abs(side_offset) > 1.0:
                    self._obstacle_side = candidate_side
                elif self._last_obstacle_side_offset is not None:
                    offset_motion = side_offset - self._last_obstacle_side_offset
                    if self._obstacle_side > 0.0 and side_offset > -3.0 and offset_motion > 2.0:
                        self._obstacle_side = -1.0
                    elif self._obstacle_side < 0.0 and side_offset < 3.0 and offset_motion < -2.0:
                        self._obstacle_side = 1.0
            self._last_obstacle_side_offset = side_offset
            self._obstacle_detection_frames += 1
            self._obstacle_missing = 0
            urgency = float(np.clip((obstacle_y - 22.0) / 18.0, 0.0, 1.0))
            steering += (
                self._obstacle_side
                * self.OBSTACLE_STEER
                * self.obstacle_steer_scale
                * urgency
            )
            target_speed = min(
                target_speed,
                self.obstacle_far_speed
                if obstacle_y < 44.0
                else self.obstacle_near_speed,
            )
        elif self._obstacle_side != 0.0:
            self._obstacle_missing += 1
            if self._obstacle_missing <= 1:
                target_speed = min(target_speed, 43.0)
            else:
                self._obstacle_side = 0.0
                self._obstacle_missing = 0
                self._last_obstacle_side_offset = None

        if self._target_speed is not None:
            target_speed = 0.65 * self._target_speed + 0.35 * target_speed
        self._target_speed = target_speed

        speed = self._observation_speed(observation, frame)
        self._last_speed_estimate = speed
        self._speed_samples += 1
        self._speed_sum += speed
        self._speed_max = max(self._speed_max, speed)
        gas, brake = self._pedals(speed, target_speed)
        self._last_target_speed = target_speed
        self._last_gas = gas
        self._last_brake = brake
        self._gas_frames += int(gas > 0.0)
        self._brake_frames += int(brake > 0.0)
        return np.asarray(
            [
                np.clip(steering, -self.MAX_STEER, self.MAX_STEER),
                gas,
                brake,
            ],
            dtype=np.float32,
        )

    def diagnostics(self) -> dict[str, float | int]:
        """Expose pixel-controller counters for evaluation without simulator state."""
        return {
            "obstacle_encounters": self._obstacle_encounters,
            "obstacle_detection_frames": self._obstacle_detection_frames,
            "pixel_speed_mean": (
                self._speed_sum / self._speed_samples if self._speed_samples else 0.0
            ),
            "pixel_speed_max": self._speed_max,
            "gas_frames": self._gas_frames,
            "brake_frames": self._brake_frames,
        }

    def last_step_diagnostics(self) -> dict[str, float | None]:
        """Expose the last pixel-based control decision for offline trace analysis."""
        obstacle = self._last_obstacle
        return {
            "target_speed": self._last_target_speed,
            "pixel_speed": self._last_speed_estimate,
            "gas": self._last_gas,
            "brake": self._last_brake,
            "obstacle_y": obstacle[0] if obstacle is not None else None,
            "obstacle_x": obstacle[1] if obstacle is not None else None,
            "road_center_at_obstacle": obstacle[2] if obstacle is not None else None,
            "obstacle_steer_side": self._obstacle_side or None,
        }


__all__ = ["VisionCorridorAgent"]

\end{lstlisting}

\clearpage
\section{haic\_agent/pixel\_features.py 전체 소스}
{\footnotesize SHA-256:\par\ttfamily f68895b9c0297ac5075ea47ab7f2a868e739b0ad24bee980082ce818b661d22a}
\begin{lstlisting}
"""Deterministic perception features derived only from the 84x84 pixel view.

These values describe HUD and road objects. They never select or modify an
action; the learned policy remains responsible for driving decisions.
"""

from __future__ import annotations

import numpy as np


ROAD_LOW = 0.24
ROAD_HIGH = 0.52
OBSTACLE_LOW = 0.54
SPEED_ROI = (77, 83, 10, 13)
SPEED_BASELINE = 0.27
SPEED_PER_UNIT = 0.085
FEATURE_NAMES = (
    "speed_fraction",
    "near_road_center_offset",
    "far_road_center_offset",
    "road_curve_magnitude",
    "obstacle_present",
    "obstacle_lateral_offset",
    "obstacle_urgency",
)
VISUAL_FEATURE_SIZE = len(FEATURE_NAMES)
TEMPORAL_FEATURE_NAMES = ("previous_obstacle_lateral_offset",)
TEMPORAL_FEATURE_SIZE = len(TEMPORAL_FEATURE_NAMES)


def current_frame(observation: np.ndarray) -> np.ndarray:
    """Validate an observation and return its newest normalized grayscale frame."""
    try:
        pixels = np.asarray(observation, dtype=np.float32)
    except (TypeError, ValueError) as error:
        raise ValueError("observation must contain normalized float pixels") from error
    if pixels.shape == (4, 84, 84):
        frame = pixels[-1]
    elif pixels.shape == (84, 84):
        frame = pixels
    else:
        raise ValueError("observation must have shape (84, 84) or (4, 84, 84)")
    if not np.all(np.isfinite(frame)) or float(frame.min()) < 0.0 or float(frame.max()) > 1.0:
        raise ValueError("observation pixels must be finite and normalized to [0, 1]")
    return frame


def road_centers(frame: np.ndarray) -> dict[int, float]:
    """Track the visible asphalt center from near to far image rows."""
    asphalt = (frame >= ROAD_LOW) & (frame <= ROAD_HIGH)
    horizontal = np.arange(frame.shape[1], dtype=np.float32)
    centers: dict[int, float] = {}
    previous = 42.0
    for row in (54, 50, 46, 42, 38, 34, 30):
        selected = asphalt[row] & (np.abs(horizontal - previous) <= 17.0)
        locations = np.flatnonzero(selected)
        if len(locations) >= 4:
            previous = float(locations.mean())
            centers[row] = previous
    return centers


def center_at(row: float, centers: dict[int, float]) -> float:
    """Interpolate a road center for a fractional image row."""
    known_rows = sorted(centers)
    if not known_rows:
        return 42.0
    return float(
        np.interp(
            row,
            np.asarray(known_rows, dtype=np.float32),
            np.asarray([centers[y] for y in known_rows], dtype=np.float32),
        )
    )


def nearest_bright_object(
    frame: np.ndarray, centers: dict[int, float]
) -> tuple[float, float, float] | None:
    """Find the nearest small bright component inside the visible road corridor."""
    if len(centers) < 3:
        return None
    bright = frame >= OBSTACLE_LOW
    bright[:22, :] = False
    bright[62:, :] = False
    visited = np.zeros(bright.shape, dtype=np.bool_)
    candidates: list[tuple[float, float, float]] = []
    height, width = bright.shape

    for start_y in range(22, 62):
        for start_x in range(width):
            if not bright[start_y, start_x] or visited[start_y, start_x]:
                continue
            stack = [(start_x, start_y)]
            visited[start_y, start_x] = True
            min_x = max_x = start_x
            min_y = max_y = start_y
            sum_x = sum_y = area = 0
            while stack:
                x, y = stack.pop()
                area += 1
                sum_x += x
                sum_y += y
                min_x = min(min_x, x)
                max_x = max(max_x, x)
                min_y = min(min_y, y)
                max_y = max(max_y, y)
                for neighbor_y in range(max(22, y - 1), min(62, y + 2)):
                    for neighbor_x in range(max(0, x - 1), min(width, x + 2)):
                        if bright[neighbor_y, neighbor_x] and not visited[neighbor_y, neighbor_x]:
                            visited[neighbor_y, neighbor_x] = True
                            stack.append((neighbor_x, neighbor_y))

            box_width = max_x - min_x + 1
            box_height = max_y - min_y + 1
            if not (4 <= area <= 80 and 2 <= box_width <= 9 and 2 <= box_height <= 10):
                continue
            center_x = sum_x / area
            center_y = sum_y / area
            road_center = center_at(center_y, centers)
            if abs(center_x - road_center) <= 12.0:
                candidates.append((center_y, center_x, road_center))

    return max(candidates, key=lambda item: item[0]) if candidates else None


def estimate_speed(frame: np.ndarray) -> float:
    """Decode the rendered speed bar using its calibrated grayscale mass."""
    top, bottom, left, right = SPEED_ROI
    mass = float(np.asarray(frame, dtype=np.float32)[top:bottom, left:right].sum())
    return float(np.clip((mass - SPEED_BASELINE) / SPEED_PER_UNIT, 0.0, 80.0))


def estimate_observation_speed(observation: np.ndarray, frame: np.ndarray) -> float:
    """Average the last two rendered speed readings in the frame stack."""
    pixels = np.asarray(observation, dtype=np.float32)
    stack = pixels[-2:] if pixels.shape == (4, 84, 84) else pixels[np.newaxis, ...]
    return float(np.mean([estimate_speed(item) for item in stack]))


def road_sweep(centers: dict[int, float]) -> float:
    """Estimate bend strength from visible near/far road-center changes."""
    return max(
        (
            abs(centers.get(row, 42.0) - centers.get(row + 24, 42.0))
            for row in (30, 34, 38, 42)
        ),
        default=0.0,
    )


def extract_visual_features(observation: np.ndarray) -> np.ndarray:
    """Encode HUD speed, road alignment, bend, and nearest obstacle from pixels."""
    frame = current_frame(observation)
    centers = road_centers(frame)
    obstacle = nearest_bright_object(frame, centers)
    if obstacle is None:
        obstacle_present = 0.0
        obstacle_lateral = 0.0
        obstacle_urgency = 0.0
    else:
        obstacle_y, obstacle_x, obstacle_road_center = obstacle
        obstacle_present = 1.0
        obstacle_lateral = float(
            np.clip((obstacle_x - obstacle_road_center) / 12.0, -1.0, 1.0)
        )
        obstacle_urgency = float(np.clip((obstacle_y - 22.0) / 18.0, 0.0, 1.0))
    return np.asarray(
        (
            estimate_observation_speed(observation, frame) / 80.0,
            np.clip((centers.get(54, 42.0) - 42.0) / 42.0, -1.0, 1.0),
            np.clip((centers.get(42, 42.0) - 42.0) / 42.0, -1.0, 1.0),
            np.clip(road_sweep(centers) / 42.0, 0.0, 1.0),
            obstacle_present,
            obstacle_lateral,
            obstacle_urgency,
        ),
        dtype=np.float32,
    )


def extract_visual_features_batch(observations: np.ndarray) -> np.ndarray:
    """Extract a row per observation for a batch of ``(B, 4, 84, 84)`` pixels."""
    pixels = np.asarray(observations, dtype=np.float32)
    if pixels.ndim != 4 or pixels.shape[1:] != (4, 84, 84):
        raise ValueError("observations must have shape (B, 4, 84, 84)")
    if not np.all(np.isfinite(pixels)) or float(pixels.min()) < 0.0 or float(pixels.max()) > 1.0:
        raise ValueError("observation pixels must be finite and normalized to [0, 1]")
    return np.stack([extract_visual_features(item) for item in pixels], axis=0)


def extract_temporal_features(observation: np.ndarray) -> np.ndarray:
    """Read a prior-frame obstacle side only when the current view is centered."""
    pixels = np.asarray(observation, dtype=np.float32)
    if pixels.shape != (4, 84, 84):
        raise ValueError("observation must have shape (4, 84, 84)")
    current = pixels[-1]
    previous = pixels[-2]
    current_centers = road_centers(current)
    current_obstacle = nearest_bright_object(current, current_centers)
    if current_obstacle is None:
        return np.zeros(TEMPORAL_FEATURE_SIZE, dtype=np.float32)
    _, current_x, current_road_center = current_obstacle
    current_lateral = float((current_x - current_road_center) / 12.0)
    if abs(current_lateral) >= 0.2:
        return np.zeros(TEMPORAL_FEATURE_SIZE, dtype=np.float32)
    previous_centers = road_centers(previous)
    previous_obstacle = nearest_bright_object(previous, previous_centers)
    if previous_obstacle is None:
        return np.zeros(TEMPORAL_FEATURE_SIZE, dtype=np.float32)
    _, previous_x, previous_road_center = previous_obstacle
    previous_lateral = float(
        np.clip((previous_x - previous_road_center) / 12.0, -1.0, 1.0)
    )
    if abs(previous_lateral) < 0.3:
        return np.zeros(TEMPORAL_FEATURE_SIZE, dtype=np.float32)
    return np.asarray((previous_lateral,), dtype=np.float32)


def extract_temporal_features_batch(observations: np.ndarray) -> np.ndarray:
    """Extract one temporal side feature per ``(B, 4, 84, 84)`` observation."""
    pixels = np.asarray(observations, dtype=np.float32)
    if pixels.ndim != 4 or pixels.shape[1:] != (4, 84, 84):
        raise ValueError("observations must have shape (B, 4, 84, 84)")
    if not np.all(np.isfinite(pixels)) or float(pixels.min()) < 0.0 or float(pixels.max()) > 1.0:
        raise ValueError("observations pixels must be finite and normalized to [0, 1]")
    return np.stack([extract_temporal_features(item) for item in pixels], axis=0)


__all__ = [
    "FEATURE_NAMES",
    "TEMPORAL_FEATURE_NAMES",
    "TEMPORAL_FEATURE_SIZE",
    "VISUAL_FEATURE_SIZE",
    "center_at",
    "current_frame",
    "estimate_observation_speed",
    "estimate_speed",
    "extract_visual_features",
    "extract_visual_features_batch",
    "extract_temporal_features",
    "extract_temporal_features_batch",
    "nearest_bright_object",
    "road_centers",
    "road_sweep",
]

\end{lstlisting}

\clearpage
\section{haic\_agent/fixed\_high\_speed\_recovery\_v2.py 전체 소스}
{\footnotesize SHA-256:\par\ttfamily dd1da6da44e009043625af401f5d1cf1cd09b915608b7fdbb42a2a8f2b46e0fc}
\begin{lstlisting}
"""Stateful reacquisition and stale-avoidance clearing at the same high target."""
import numpy as np
from haic_agent.fixed_high_speed_runtime import FixedHighSpeedAgent

MODES = ("memory_reacquire", "impact_clear")


class FixedHighSpeedRecoveryV2(FixedHighSpeedAgent):
    def __init__(self, mechanism="memory_reacquire"):
        if mechanism not in MODES:
            raise ValueError("unknown recovery mode")
        self.recovery_mode = mechanism
        super().__init__("damping")

    def reset(self, observation=None):
        super().reset(observation)
        self.remembered_road = 0.0
        self.latched_turn = 0.0
        self.recovery_left = 0
        self.recovery_armed = True
        self.centered_streak = 0
        self.impact_left = 0
        self.speed_history = []

    def act(self, observation):
        action = super().act(observation)
        original = float(action[0])
        centers = self.last["road_centers"]
        speed = self.last["pixel_speed"]
        complete = 42 in centers and 54 in centers
        centered = len(centers)==7 and abs(centers[42]-42)<=8 and abs(centers[54]-42)<=8
        self.centered_streak = self.centered_streak+1 if centered else 0
        if self.centered_streak >= 3:
            self.recovery_left = 0
            self.recovery_armed = True
        trigger = False
        if (self.steps>10 and self.impact_left==0 and self.speed_history and speed<50
                and max(self.speed_history)-speed>10):
            self.impact_left = 12
            trigger = True
        if self.steps>10 and self.recovery_mode=="memory_reacquire":
            if (42 not in centers and self.recovery_armed and abs(self.remembered_road)>.01):
                self.recovery_left = 80
                self.recovery_armed = False
                self.latched_turn = np.sign(self.remembered_road)*np.clip(abs(self.remembered_road),.35,.7)
            if self.recovery_left>0:
                action[0] = self.latched_turn
                self.recovery_left -= 1
        elif self.steps>10 and self.recovery_mode=="impact_clear" and self.impact_left>0 and complete:
            road = .022*(centers[42]-42)+.018*(centers[42]-centers[54])
            action[0] = np.clip(road+self.last["correction"],-.7,.7)
        if complete and self.recovery_left==0:
            self.remembered_road = .022*(centers[42]-42)+.018*(centers[42]-centers[54])
        if self.impact_left:
            self.impact_left -= 1
        self.speed_history = (self.speed_history+[speed])[-4:]
        self.last.update(recovery_mode=self.recovery_mode,damping_steer=original,
                         recovery_changed=abs(float(action[0])-original)>1e-6,
                         recovery_steps_remaining=self.recovery_left,centered_streak=self.centered_streak,
                         impact_proxy_trigger=trigger,impact_steps_remaining=self.impact_left)
        return action

\end{lstlisting}

\clearpage
\section{haic\_agent/fast\_completion\_coordination.py 전체 소스}
{\footnotesize SHA-256:\par\ttfamily 2f890619535b744977ecbee04c1604f9bed8f2be94fd92de733dfb1c0c895d84}
\begin{lstlisting}
"""Coordinate high-speed pedals and steering using only rendered road cues."""
import numpy as np
from haic_agent.fixed_high_speed_recovery_v2 import FixedHighSpeedRecoveryV2
from haic_agent.pixel_features import road_centers

MODES=('control','preview_budget','actuator_rate','impact_traction','recovery_budget',
       'preview_action_aware','recovery_geometry_aware','preview_row_repair')


def observed_center(centers,row):
    if row in centers: return centers[row]
    known=sorted(centers,key=lambda y:abs(y-row))
    if len(known)<2: return None
    a,b=known[:2]
    return float(np.clip(centers[a]+(row-a)*(centers[b]-centers[a])/(b-a),0,83))


class FastCompletionCoordination(FixedHighSpeedRecoveryV2):
    def __init__(self,mechanism='control'):
        if mechanism not in MODES: raise ValueError('unknown mechanism')
        self.coordination_mode=mechanism
        super().__init__('impact_clear')

    def reset(self,observation=None):
        super().reset(observation)
        self.last_command=0.
        self.recovering=False
        self.stable=0
        self.brake_history=[]
        self.visible_steer=0.
        self.lost_rows=0

    def act(self,observation):
        impact_before=self.impact_left
        action=super().act(observation)
        original=action.copy()
        centers=self.last['road_centers']
        speed=self.last['pixel_speed']
        target=60.
        veto=False
        geometry_repaired=False
        if self.steps>10:
            if self.coordination_mode in ('preview_action_aware','recovery_geometry_aware','preview_row_repair'):
                veto=self.last['impact_proxy_trigger'] and any(b>.01 for b in self.brake_history)
                if veto:
                    action[0]=self.last['damping_steer']
                    self.impact_left=0
            if self.coordination_mode=='preview_row_repair' and 42 not in centers and 54 in centers and len(centers)>=2:
                m=observed_center(centers,42)
                obstacle=self.base._last_obstacle
                term=self.base._obstacle_side*.34*float(np.clip((obstacle[0]-22)/18,0,1)) if obstacle is not None else 0.
                action[0]=np.clip(.022*(m-42)+.018*(m-centers[54])+term+self.last['correction'],-.7,.7)
                geometry_repaired=True
            if self.coordination_mode=='recovery_geometry_aware':
                if 42 not in centers or 54 not in centers:
                    n,m=observed_center(centers,54),observed_center(centers,42)
                    if n is not None and m is not None:
                        stack=np.asarray(observation)
                        prev=road_centers(stack[-2]) if stack.ndim==3 else centers
                        pm=observed_center(prev,42)
                        obstacle=self.base._last_obstacle
                        clearing=(impact_before>0 or self.last['impact_proxy_trigger']) and not veto
                        term=self.base._obstacle_side*.34*float(np.clip((obstacle[0]-22)/18,0,1)) if obstacle is not None and not clearing else 0.
                        action[0]=np.clip(.022*(m-42)+.018*(m-n)+term+(.025*(m-pm) if pm is not None else 0.),-.7,.7)
                        self.visible_steer=float(action[0]);self.lost_rows=0
                        geometry_repaired=True
                    else:
                        self.lost_rows+=1
                        if self.lost_rows<=12:
                            action[0]=self.visible_steer
                            geometry_repaired=True
                else:
                    self.visible_steer=float(action[0]);self.lost_rows=0
            if self.coordination_mode in ('preview_budget','preview_action_aware','preview_row_repair'):
                sweep=max(centers.values())-min(centers.values()) if len(centers)>1 else 28.
                target=max(38.,60.-.8*sweep)
                if self.base._last_obstacle is not None: target=min(target,44.)
            elif self.coordination_mode=='actuator_rate':
                action[0]=np.clip(action[0],self.last_command-.24,self.last_command+.24)
            elif self.coordination_mode=='impact_traction':
                if impact_before>0 or self.last['impact_proxy_trigger']:
                    action[1]=min(action[1],.18)
            elif self.coordination_mode in ('recovery_budget','recovery_geometry_aware'):
                error=abs(centers.get(54,0)-42)
                if error>6 or 42 not in centers: self.recovering=True
                self.stable=self.stable+1 if len(centers)==7 and error<4 else 0
                if self.stable>=3: self.recovering=False
                if self.recovering: target=40.
            if target<60:
                action[1]=np.clip(.12+.04*(target-speed),0,.6)
                action[2]=np.clip(.02*(speed-target-2),0,.28)
        self.last_command=float(action[0])
        self.brake_history=(self.brake_history+[float(action[2])])[-4:]
        self.last.update(coordination_mode=self.coordination_mode,target_speed=target,
                         braking_proxy_veto=bool(veto),geometry_repaired=geometry_repaired,
                         completion_changed=bool(np.max(np.abs(action-original))>1e-6),
                         final_gas=float(action[1]),final_brake=float(action[2]))
        return action

\end{lstlisting}
\end{document}
